% Unofficial University of Bridgeport Poster Template
% https://github.com/andiac/gemini-cam
% a fork of https://github.com/anishathalye/gemini

\documentclass[final]{beamer}

% ====================
% Packages
% ====================

\usepackage[T1]{fontenc}
\usepackage{lmodern}
% \usepackage[size=custom,width=120,height=72,scale=1.0]{beamerposter}
\usepackage[size=custom,width=91.44,height=60.96,scale=1.04] {beamerposter} % centimeter - Standard Landscape

%\usepackage[size=custom,width=60.96,height=91.44,scale=1.0] {beamerposter} % centimeter - UBRise Portrait

\usetheme{gemini}
\usecolortheme{UBridgeport}
\usepackage{graphicx}
\usepackage{booktabs}
\usepackage{tikz}
\usepackage{pgfplots}
\usepgfplotslibrary{colorbrewer}
\usepgfplotslibrary{statistics}
\pgfplotsset{compat=1.18}
\usepackage{pgf-pie}
\usepackage{anyfontsize}

% ====================
% Lengths
% ====================

% If you have N columns, choose \sepwidth and \colwidth such that
% (N+1)*\sepwidth + N*\colwidth = \paperwidth
\newlength{\sepwidth}
\newlength{\colwidth}
\setlength{\sepwidth}{0.025\paperwidth}
\setlength{\colwidth}{0.3\paperwidth}

\newcommand{\separatorcolumn}{\begin{column}{\sepwidth}\end{column}}

% ====================
% Title
% ====================
\section{Header}
\title{Cleaning up Seaside Park: A brief analysis on beach pollution along Long Island Sound}

\author{Huy Huong\inst{1} \and Samantha Targonski\inst{1} \and Alexandra Mones\inst{2} \and Emmanuel Ugbomah\inst{2}}

\institute[UB]{HONR-390/CAPS-390 Senior Capstone --- Island of Misfits} 

% ====================
% Footer (optional)
% ====================
\section{Footer}
\footercontent{
  { } \hfill
  HONR-390/CAPS-390 Senior Capstone --- Island of Misfits \hfill
  \href{mailto:john.doe@example.com}{ }}
% (can be left out to remove footer)s

% ====================
% Logo (optional)
% Refer to https://github.com/k4rtik/uchicago-poster
% logo: https://communications.admin.ox.ac.uk/communications-resources/visual-identity/identity-guidelines/the-oxford-logo
% use this to include logos on the left and/or right side of the header:
% ====================

\logoleft{\includegraphics[height=7cm]{logos/ub-logo-stacked-KO.png}}
\logoright{\includegraphics[height=7cm]{logos/ub-logo-stacked-KO.png}}

% ====================
% Body
% ====================

\begin{document}
  \begin{frame}[t]
    \begin{columns}[t]
      \separatorcolumn

      \begin{column}{\colwidth}
        \section{Abstract}
        \begin{block}{Introduction}

          Long Island Sound, a vital estuary located between Connecticut and Long Island, serves as an ecological and economic hub, supporting diverse marine life, recreational activities, and commercial industries. However, like many coastal regions, the Sound faces significant environmental challenges, with one of the most pressing issues being the growing problem of litter. From plastic debris and discarded fishing gear to food wrappers and cigarette butts, litter in the Sound poses threats to water quality, wildlife, and human health. The accumulation of waste not only harms marine species but also affects local economies that rely on tourism and fishing. While organizations like Save the Sound \cite{STS_CT_Cleanup} work across Connecticut to provide resources and organize local cleanups to reduce overall pollution, more action is needed. This research poster examines the sources, types, and consequences of litter in Long Island Sound and explores potential solutions to mitigate this issue, protecting this invaluable resource for future generations.

        \end{block}

        \begin{block}{Research Objectives}

          This research examines the volume of litter collected at Seaside Park on a single day to estimate annual waste generation and identify local sources of pollution. The study aims to also determine the most common types of litter found in the park and gather survey data to develop strategies for reducing waste. With both quantitative and qualitative data, we hope to contribute to broader efforts in environmental sustainability.

        \end{block}

        \begin{block}{Methodology}
            \begin{itemize}
              \item A small survey taken by the Fall 2024 Senior Capstone class on beach pollution
              \item Collected field research by going to Seaside Park and recording litter found along a section of the walkway.
            \end{itemize}

        \end{block}
      \end{column}

      \separatorcolumn

      \begin{column}{\colwidth}

        \begin{block}{Results}
            \begin{figure}[h!]
                \centering
                % First Pie Chart
                \begin{minipage}[t]{0.45\textwidth}
                    \centering
                    \begin{tikzpicture}
                      \pie[
                        radius=3.5,
                        text=legend,
                        font=\small,
                        color={purple, pink, orange, green}
                    ]{
                        7.7/Always,
                        61.5/Often,
                        23.1/Sometimes,
                        7.7/Rarely
                    }
                \end{tikzpicture}
                \caption{\centering How often do you see trash at CT beaches?}
                \end{minipage}%
                \hfill
                % Second Pie Chart
                \begin{minipage}[t]{0.45\textwidth}
                    \centering
                    \begin{tikzpicture}
                        \pie[
                            radius=3.5,
                            text=legend,
                            font=\small,
                            color={purple, 
                                   pink, 
                                   orange, 
                                   green, 
                                   cyan}
                        ]{
                            7.7/Always,
                            7.7/Often,
                            15.4/Sometimes,
                            30.8/Rarely,
                            38.5/Never
                        }
                    \end{tikzpicture}
                    \caption{\centering How often do you pick up litter when you visit a beach?}
                \end{minipage}
            \end{figure}
            Our survey, conducted via a Google Form, garnered 13 responses addressing beach pollution. While students demonstrated awareness of the issue, their actions tell a different story. Notably, 9 out of 13 respondents indicated that they Rarely or Never pick up litter during beach visits. Health concerns were a commonly cited reason for avoiding direct engagement in cleanup efforts.
            \\
            The challenge of addressing beach pollution is not just about awareness but also the associated costs. Researchers from the Center for Marine Socioecology at the University of Tasmania estimate that achieving just a 1\% reduction in global pollution would require an annual contribution of US \$0.71 per person worldwide \cite{MUTUKU2022114220}.
          \begin{table}
            \centering
            \begin{tabular}{l r }
              \toprule
              \textbf{Item} & \textbf{Quantity} \\
              \midrule
              Paper \& Plastic Pieces& 87\\
              Bottles/Cups/Cans& 63\\
              Glass& 15\\
              Takeout Containers& 10\\
              Bags& 8\\
              Large Items& 4\\
              \bottomrule
            \end{tabular}
            \caption{Litter picked up on 11/25/24.}
          \end{table}
          During our walk around Seaside Park, we observed a significant amount of litter throughout the area. Among the 187 items collected, plastic and paper pieces accounted for 46\% of the total. Additionally, we discovered numerous plastic bottles, cups, and cans scattered and hidden behind the rock wall along the beach. Many of these items displayed clear signs of oxidation and deterioration, indicating they had been left uncollected for an extended period.
        \end{block}

      \end{column}

      \separatorcolumn

      \begin{column}{\colwidth}

        \begin{block}{Discussion}
        
          Use this section to discuss the results and how to interpret them \par
          Class aptent taciti sociosqu ad litora torquent per conubia nostra, per
          inceptos himenaeos. Phasellus libero enim, gravida sed erat sit amet,
          scelerisque congue diam. Fusce dapibus dui ut augue pulvinar iaculis.

        \end{block}

        \begin{block}{Conclusion}
        
          Can someone write here a short conclusion on the entire project and then we can edit from there? \par
          Class aptent taciti sociosqu ad litora torquent per conubia nostra, per
          inceptos himenaeos. Phasellus libero enim, gravida sed erat sit amet,
          scelerisque congue diam. Fusce dapibus dui ut augue pulvinar iaculis.

        \end{block}

        \begin{block}{References}

          \nocite{*}
          \footnotesize{\bibliographystyle{plain}\bibliography{poster}}

        \end{block}  
      \end{column}

      \separatorcolumn
    \end{columns}
  \end{frame}

\end{document}